\section{Introduction}\label{sec:intro}

Scheduled macroeconomic releases are the cleanest natural experiments in finance: the
time of the shock is known to the second, its content is unknown until release, and
its surprise relative to expectations can be measured. How fast prices absorb such
information, and how much of the post-release activity is information versus the
market reacting to itself, are old questions
\citep{andersen2003micro,rambaldi2015modeling}. Hawkes processes answer the second
question elegantly---the branching structure separates exogenous (news-driven)
events from endogenous (self-excited) ones---and Rambaldi, Pennesi and Lillo
\citep{rambaldi2015modeling} showed that adding a news kernel to a Hawkes model of FX
quote activity captures the post-announcement burst.

Three gaps remain. \emph{First}, the absorption speed is usually read off the news
kernel alone (the ``half-life'' $\ln 2/\beta$ of an exponential), which ignores the
endogenous echoes triggered by the news; when the branching ratio is close to one
\citep{hardiman2013critical,filimonov2012quantifying} the echoes dominate the total
response. \emph{Second}, estimation is fragile: flexible kernels are fitted by local
optimisation of a non-convex likelihood, and branching ratios are biased upward when
seasonality is mis-specified \citep{filimonov2015apparent}. \emph{Third}, the
information content of the release (its surprise) enters, if at all, as a post-hoc
correlate rather than as part of the model, and cross-market propagation is rarely
modelled jointly.

\paragraph{Contributions.} Linear-in-parameters kernel dictionaries, EM for Hawkes
processes and the Markov embedding of Erlang-kernel Hawkes processes are known
individually; our contribution is to combine them into a measurement framework whose
output---echo-corrected absorption times---is certified and falsifiable.
\begin{enumerate}
\item \textbf{A certified estimator.} We represent every kernel---endogenous, news,
anticipatory---in a phase-type dictionary (exponential and Erlang phases on a log grid
of time scales). All parameters enter the intensity linearly, the log-likelihood is
concave, and we compute the global maximum with accelerated EM and a barrier-Newton
polish, returning a Fenchel duality gap that certifies optimality
(Proposition~\ref{prop:gap}).
\item \textbf{Echo-corrected absorption in closed form.} A Markov embedding turns the
mean response to a release into $Ce^{\mathcal At}\bm z_0$. We prove that the embedding is
stable iff the spectral radius of the branching matrix is below one
(Theorem~\ref{thm:stab}) and obtain absorption quantiles, echo amplification and---via
$\delta$-crossing sampling---exact expected price-discovery curves
(Corollary~\ref{cor:abs}).
\item \textbf{Identification by design.} Per-window baselines, time-of-day hats
identified by matched placebo windows, and a PLACEBO news type whose fitted kernel must
vanish give a falsification test built into the likelihood.
\item \textbf{Surprise as a mark.} Standardised surprises from real-time nowcasts
(Cleveland Fed CPI nowcast, Atlanta Fed GDPNow) and real-time statistical expectations
enter the news kernel linearly and separately by sign.
\item \textbf{EPT-TPP.} A neural temporal point process whose latent state is a
nonnegative phase-type system: exact compensator, Bayes-optimal predictions, and it
nests both multivariate Hawkes processes and RMTPP (Section~\ref{sec:ept}).
\item \textbf{Benchmarks.} Synthetic recovery against six classical estimators, public
neural-TPP benchmarks under the S2P2 protocol, LOBSTER limit-order-book streams, and
post-release forecasting on a 2022--2026 six-asset news panel.
\end{enumerate}
